\section{Conclusion and Future Directions}
\label{sec:conclusion}

In this work, we have explored an alternative to the foundational, linear paradigm of multi-task model merging. 
By introducing \textbf{FoldMerge (Neural Origami)}, we demonstrate that model merging can be formulated as an exploratory non-linear coordinate transformation using continuous weight-space diffeomorphisms rather than flat linear interpolations. 
By training a differentiable Normalizing Flow (RealNVP) with an implicit parameter-wise $\ell_2$ weight-decay flow regularization penalty, FoldMerge warps the underlying weight-coordinate system to merge separate expert basins of attraction in Origami Space. 

Our empirical evaluations on the 8-task Vision-Language ViT-B/32 benchmark provide an encouraging proof-of-concept for the viability and trainability of this approach. 
FoldMerge achieves an Average Accuracy of \textbf{89.76\%}, performing on par with the highly competitive, state-of-the-art SyMerge framework on average while outperforming it slightly on 5 out of 8 individual tasks. 
Furthermore, once optimized and decoded back, FoldMerge introduces zero inference or parameter overhead during actual model deployment, making it a viable and attractive framework for further development.

\paragraph{Future Directions:}
We believe that FoldMerge opens up promising research avenues to address the limitations identified in this work:
\begin{enumerate}
  \item \textbf{Permutation-Equivariant Flow Architectures:} To overcome the coordinate-dependence of RealNVP affine coupling layers, a crucial future direction is to design permutation-equivariant normalizing flow architectures that naturally respect neural network weight symmetries.
  \item \textbf{Unified Tensor-Aware Warping:} Rather than relying on the row-wise slicing heuristic used in this work (which introduces a structural category error by treating weight matrices as independent row-wise samples), future research should investigate tensor-aware coordinate warps that operate on entire layers or low-rank adapters as unified algebraic structures.
  \item \textbf{Parameter and Computation Compression:} We aim to explore weight-sharing and low-rank parameterizations of the scale and translation networks within the flow coupling layers to drastically reduce the flow's parameter footprint ($\approx 2.6\text{M}$ parameters) and accelerate test-time optimization.
  \item \textbf{Theoretical Foundations in Differential Topology:} We hope this work inspires deeper theoretical investigations bridging differential geometry, algebraic topology, and weight-space optimization landscapes, ultimately helping the community build a more complete, mathematically rigorous understanding of deep weight-space geometry.
\end{enumerate}
Our work marks an exploratory step towards a non-linear geometric understanding of deep weight spaces, demonstrating that non-linear parameter-space coordinate warping is a viable, computationally trainable pathway for deep model alignment and fusion.
